% Propietario conservado: /Users/ruben/Documents/ChatGPT/jueces y controles/REVISION_ARGUMENTAL_APERTURAS_CIERRES_20260915/II/manuscript_es/sections/10d_portador_pentadico_rev06.tex
% RESULTADO_RECUPERADO. Integral activo, capítulos 63--65.
% B023_ch63_descenso_icosaedrico_portador_6ba97d5bb029_07_gravedad.tex:196--380.
% B024_ch64_entrelazador_tensorial_reduccion_6ba97d5bb029_07_gravedad.tex:1--122.
% B025_ch65_descomposicion_helicoidal_espin_dos_6ba97d5bb029_07_gravedad.tex:1--104.
% Adecuación editorial: jerarquía, referencias locales y distinción del proyector icosaédrico.
\subsection{Incidence icosaédrique et couplage gravitationnel à cinq composantes}
\label{ii:sec:ii-pentadica}

La structure discrète du continu conserve une lecture d’incidence
à côté de la lecture arithmétique. Le porteur gravitationnel à cinq composantes se
construit au moyen de cette incidence et de sa représentation tensorielle. La section
de couplage \(G_a(\vartheta^\circ_{108})\), avec
\(a\in\{\mathrm{pre},\mathrm{ret}\}\), est celle obtenue dans la
section~\ref{ii:sec:gravity} ; on explicite ci-après l’espace sur
lequel elle agit, ses cinq composantes et sa projection transversale.

Soit \(X_{\rm TPK}^{\rm exc}\) le sous-espace de l’état \(\TPK\) enrichi
qui conserve, outre la sortie arithmétique, l’orientation, la route,
la retenue et l’incidence exceptionnelle. Pour expliciter l’espace géométrique d’arrivée,
soit \(\varphi_g=(1+\sqrt5)/2\), \(\nu=\sqrt{1+\varphi_g^2}\), et définissons
\begin{equation}
\label{ii:ii:pent:eq:gravity-icosahedron-vertices}
\Omega_{12}^{\rm ico}
=\frac1\nu\bigl(
\{0\}\times\{\pm1\}\times\{\pm\varphi_g\}
\;\cup\;
\{\pm1\}\times\{\pm\varphi_g\}\times\{0\}
\;\cup\;
\{\pm\varphi_g\}\times\{0\}\times\{\pm1\}
\bigr)\subset S^2.
\end{equation}
Ses douze éléments sont les classes orientées de sommet d’un icosaèdre.
L’incidence ne reste pas verbale : elle est fixée par
\begin{equation}
\label{ii:ii:pent:eq:gravity-icosahedron-incidence}
\mathcal I_{12}^{\rm ico}
=\{(v,w)\in\Omega_{12}^{\rm ico}\times\Omega_{12}^{\rm ico}:
v\neq w,\ \langle v,w\rangle=1/\sqrt5\}.
\end{equation}
Chaque sommet a cinq voisins dans cette relation. Le pont de lecture
est typé comme une application
\begin{equation}
\label{ii:ii:pent:eq:gravity-exceptional-reader-map}
\pi_{\rm ico}:X_{\rm TPK}^{\rm exc}\longrightarrow\Omega_{12}^{\rm ico}.
\end{equation}
La réalisation d’incidence exige, outre les douze éléments de
l’image, les propriétés simultanées suivantes de l’application :
\begin{enumerate}[label=(\roman*)]
\item être surjective et constante exactement sur les fibres qui écartent
la carte arithmétique mais conservent la route exceptionnelle ;
\item entrelacer l’action des automorphismes orientés de l’antécédent
avec une action transitive sur \(\Omega_{12}^{\rm ico}\) ;
\item transporter l’incidence TPK déclarée vers
\(\mathcal I_{12}^{\rm ico}\), et non
ses seuls cardinaux ;
\item conserver l’involution de feuillet et le retour nonadique dans le
stabilisateur de la fibre correspondante.
\end{enumerate}

Ces conditions fixent l’ordre de la construction : d’abord l’application
\(\pi_{\rm ico}\), puis l’identification de l’action induite et,
enfin, sa représentation linéaire. L’orientation, le chemin et la retenue
appartiennent à l’antécédent enrichi de cette descente.


% BEGIN OWNER /Users/ruben/Documents/ChatGPT/jueces y controles/REVISION_ARGUMENTAL_APERTURAS_CIERRES_20260915/II/manuscript_es/sections/representacion_dodecafasica_rev08.tex
% FORMALIZACION_REUNIDA de la proyección del calendario y la representación
% de doce posiciones ya desarrolladas en la arquitectura TPK y el enlace Witt.
% Formalización explícita del cotejo de Gram y del transporte de marco.
\subsubsection{Projection du calendrier et représentation de ses douze positions}
\label{ii:ii:dodeca:calendario}
Le calendrier de la section~\ref{ii:ii:tpk-arquitectura} fournit une
coordonnée observable de l’état complet. Avec ses notations, nous écrivons
\(r_0(\theta)=r(\theta)-1\in\{0,\ldots,8\}\) et
\(b=\beta(\theta)\in\mathbb Z/12\mathbb Z\). La projection d’horloge est
\[
 Q_{\rm clk}(x)=\bigl(\beta(\theta(x)),r_0(\theta(x))\bigr).
\]
Pour un chemin admissible de \(n\) mises à jour, le transport induit est
\begin{equation}
 k(r_0,n)=\left\lfloor\frac{r_0+n}{9}\right\rfloor,
 \qquad
 T_n(b,r_0)=\bigl(b+k(r_0,n)\bmod12,\ (r_0+n)\bmod9\bigr).
 \label{ii:ii:dodeca:transporte-reloj}
\end{equation}
Ici, \(n\in\mathbb N_0\) compte les mises à jour effectives ; le quotient
entier \(k\) enregistre les retours nonadiques traversés.

\begin{proposition}[Compatibilité du calendrier avec la composition des chemins]
\label{ii:ii:dodeca:composicion}
La projection \(Q_{\rm clk}\), avec la longueur de chaque chemin, définit
un foncteur de la catégorie des chemins admissibles TPK vers la catégorie d’action
du calendrier~\eqref{ii:ii:dodeca:transporte-reloj}. On a
\(T_{m+n}=T_mT_n\), \(T_0=I\),
\(T_9(b,r_0)=(b+1,r_0)\) et \(T_{108}(b,r_0)=(b,r_0)\).
\end{proposition}
\begin{proof}
En posant \(r'_0=(r_0+n)\bmod9\), la division euclidienne donne
\[
 k(r_0,n+m)=k(r_0,n)+k(r'_0,m).
\]
Cette identité conserve simultanément le résidu et le quotient lors de la concaténation
des chemins. La mise à jour de \(\theta\) établie dans la
section~\ref{ii:ii:tpk-arquitectura} induit exactement
\eqref{ii:ii:dodeca:transporte-reloj} ; les identités restantes s’obtiennent
par substitution. La catégorie d’arrivée a les paires \((b,r_0)\) pour
objets et les longueurs admissibles pour flèches, avec la somme comme composition.
\end{proof}

L’indice \(p=b+1\in\{1,\ldots,12\}\), en prenant le représentant
\(0\le b<12\), publie la position dodécaphasique. Un tour de
\(108\) mises à jour parcourt ses douze positions. La périodicité
démontrée correspond à cette projection : route, retenue accumulée et
mémoire restent dans l’état source et poursuivent leur mise à jour.

\paragraph{Réalisation d’incidence à partir du projecteur angulaire.}
Les représentants \(r_pC_5\) et la représentation \(\rho\) de la
section~\ref{ii:sec:ii-enlace-angular-witt} fixent la carte des douze
positions. Le projecteur déjà défini par
\eqref{ii:eq:ii-witt-projector-character} produit
\begin{equation}
 \begin{gathered}
 \mathcal H_3=P_3^{\mathrm{ico}}\mathbb R^{12},\qquad v_p=2P_3^{\mathrm{ico}}e_p,
 \\[-2pt]
 (P_3^{\mathrm{ico}})_{pq}=\frac1{20}
 \sum_{\substack{g\in A_5\\gr_qC_5=r_pC_5}}\chi_3(g^{-1}).
 \end{gathered}
 \label{ii:ii:dodeca:vertices-proyector}
\end{equation}
La dernière somme contient cinq termes par entrée et utilise les
représentants et les deux classes de cycles d’ordre cinq qui y sont fixés.

\begin{proposition}[Réalisation icosaédrique des positions dodécaphasiques]
\label{ii:ii:dodeca:realizacion}
Les vecteurs \(v_p\) sont douze vecteurs unitaires distincts de
\(\mathcal H_3\). Leur incidence est
\(p\sim q\Longleftrightarrow\langle v_p,v_q\rangle=1/\sqrt5\).
L’application \(p\mapsto v_p\) est \(A_5\)-équivariante et son image
est isométrique à \(\Omega_{12}^{\rm ico}\).
\end{proposition}
\begin{proof}
La transitivité et \(\operatorname{Tr}P_3^{\mathrm{ico}}=3\) donnent
\((P_3^{\mathrm{ico}})_{pp}=1/4\) ; par idempotence,
\(\langle v_p,v_q\rangle=4(P_3^{\mathrm{ico}})_{pq}\) et \(\|v_p\|=1\).
L’évaluation des cinq termes de
\eqref{ii:ii:dodeca:vertices-proyector} donne les paires antipodales
\[
 (1,4),\ (2,3),\ (5,8),\ (6,7),\ (9,12),\ (10,11).
\]
Pour les représentants \((1,2,5,6,9,10)\), leur matrice de Gram est
\begin{equation}
 I_6+\frac D{\sqrt5},\qquad
 D=\begin{pmatrix}
 0&-1&-1&1&-1&1\\
 -1&0&1&1&-1&-1\\
 -1&1&0&1&1&1\\
 1&1&1&0&-1&1\\
 -1&-1&1&-1&0&1\\
 1&-1&1&1&1&0
 \end{pmatrix},\qquad D^2=5I_6.
 \label{ii:ii:dodeca:gram}
\end{equation}
Les autres produits scalaires s’obtiennent par antipodie.
Ainsi, hormis un sommet et son opposé, les produits sont
\(\pm1/\sqrt5\), et chaque sommet a cinq voisins. Le tableau suivant
explicite la comparaison avec la carte signée à six positions
\(\mathcal U=(\infty,0,1,2,3,4)\) :
\begin{equation}
\begin{array}{c|cc@{\qquad}c|cc}
p&u(p)&\sigma(p)&p&u(p)&\sigma(p)\\ \hline
1&\infty&+1&7&3&-1\\
2&0&-1&8&1&+1\\
3&0&+1&9&2&-1\\
4&\infty&-1&10&4&+1\\
5&1&-1&11&4&-1\\
6&3&+1&12&2&+1
\end{array}
\label{ii:ii:dodeca:carta-firmada}
\end{equation}
En effet, pour \(C=\operatorname{bal}(A_W)\), avec \(A_W\) donné dans
\eqref{ii:exc:aw}, les entrées de \(D\) sont
\(D_{ab}=\sigma(p_a)\sigma(p_b)C_{u(p_a),u(p_b)}\).
Par conséquent, le tableau et~\eqref{ii:ii:dodeca:gram} identifient
isométriquement \(v_p\) à
\(\sigma(p)\sqrt2 E_Ce_{u(p)}\), où
\(E_C=(I_6+C/\sqrt5)/2\). Les deux systèmes engendrent des espaces de dimension
trois et ont une matrice de Gram identique, ce qui prouve que l’application
linéaire induite est bien définie et isométrique. La réalisation
explicite de~\eqref{ii:ii:ico:proyector-seis}, démontrée ci-après,
l’identifie à~\(\Omega_{12}^{\rm ico}\) et à son incidence.
Enfin, \(P_3^{\mathrm{ico}}\rho(g)=\rho(g)P_3^{\mathrm{ico}}\) implique
\(v_{g\cdot p}=\rho(g)v_p\). La bijection des douze positions donne
\(\operatorname{Stab}_{A_5}(v_p)=r_pC_5r_p^{-1}\).
\end{proof}

On obtient ainsi la composition représentationnelle
\begin{equation}
 x\longmapsto Q_{\rm clk}(x)\longmapsto
 p=\beta(\theta(x))+1\longmapsto v_p.
 \label{ii:ii:dodeca:composicion-representacional}
\end{equation}
L’expression signée de cette composition est
\(r_\pm^{\rm rep}(x)=(u(p),\sigma(p))\), déterminée par le
tableau~\eqref{ii:ii:dodeca:carta-firmada}. Le tableau constitue un choix explicite
d’alignement entre cartes ; les actions de \(A_5\) fournissent ses
alignements équivalents. Le registre complet est conservé dans le
domaine de~\eqref{ii:ii:dodeca:composicion-representacional}.

\paragraph{Lien avec le porteur tensoriel à cinq composantes.}
Dans \(\mathcal H_3\), nous formons
\(Q_p=v_pv_p^*-I_{\mathcal H_3}/3\). Soit \(J=\mathbf1\mathbf1^*\)
et soit \(\mathsf A\) la matrice de la permutation antipodale précédente.
Le produit de Frobenius et~\eqref{ii:ii:dodeca:gram} donnent
\begin{equation}
 \begin{aligned}
 \operatorname{Gram}(Q_p)
 &=(4P_3^{\mathrm{ico}})\circ(4P_3^{\mathrm{ico}})-\frac13J
 =\frac45(I+\mathsf A)-\frac2{15}J,
 \\
 P^{\rm ico}_5&=\frac58\operatorname{Gram}(Q_p)
 =\frac12(I+\mathsf A)-\frac1{12}J.
 \end{aligned}
 \label{ii:ii:dodeca:puente-tensorial}
\end{equation}
Le symbole \(\circ\) désigne le produit entrée par entrée. La dernière
expression est le projecteur orthogonal sur les coordonnées paires sous
antipodie et orthogonales à \(\mathbf1\), de dimension cinq. Son caractère
est \((5,1,-1,0,0)\) : sur les six paires antipodales, les nombres de
points fixes sont \((6,2,0,1,1)\), et l’on retire la composante constante.
Il coïncide donc avec le projecteur central de
\eqref{ii:ii:pent:eq:gravity-p5}. Cette identité relie la construction
tridimensionnelle à \(\Gamma_0\) et \(\Gamma_5\), définis dans
\eqref{ii:ii:pent:eq:gravity-Gamma0} et~\eqref{ii:ii:pent:eq:gravity-gamma5},
au moyen de l’opération quadratique explicite \(v_p\mapsto Q_p\).

\subsubsection{Transport du repère représentationnel}
\label{ii:ii:dodeca:marco}
Soit \(s=(1\,2\,\cdots\,12)\) et soit \(S_ce_p=e_{s(p)}\).
Un retour nonadique du calendrier envoie \(p\) sur \(s(p)\). Le transport
conjoint du repère, de la représentation et du projecteur s’exprime par
\begin{equation}
 \begin{gathered}
 P_{3,k}^{\mathrm{ico}}=S_c^kP_3^{\mathrm{ico}}S_c^{-k},\qquad
 \rho^{(k)}(g)=S_c^k\rho(g)S_c^{-k},\\
 v^{(k)}_{s^k(p)}=2P_{3,k}^{\mathrm{ico}}e_{s^k(p)}=S_c^kv_p.
 \end{gathered}
 \label{ii:ii:dodeca:marco-movil}
\end{equation}
La dernière égalité se prouve par substitution ; l’orthogonalité de
\(S_c\) conserve tous les produits scalaires et transporte
l’incidence. Pour un chemin de longueur \(n\), on prend
\(k=k(r_0,n)\). L’identité de retenue de la
proposition~\ref{ii:ii:dodeca:composicion} prouve la compatibilité de
ces transports avec la concaténation. Bien que \(S_c^{12}=I\),
les douze tours mettent à jour la mémoire du chemin source.

L’égalité~\eqref{ii:ii:dodeca:marco-movil} établit la covariance entre
repères. L’action active sur une incidence fixe a un contenu
supplémentaire : elle exige d’identifier le transport de l’état complet et son
action dans cette carte. En particulier, \(s\) est d’ordre douze, tandis
que les ordres des éléments de \(A_5\) sont \(1,2,3,5\).
De même, le stabilisateur d’un sommet est d’ordre cinq, et celui de sa
paire antipodale est d’ordre dix ; l’involution de la carte signée
échange les deux sommets. Ces distinctions conservent le sens
de l’orientation, du feuillet et du retour dans les propriétés (i)--(iv) : la
réalisation représentationnelle démontrée et la descente de la route
exceptionnelle complète sont reliées par ces propriétés, avec leurs
domaines et fibres déclarés.

% END OWNER


% BEGIN OWNER /Users/ruben/Documents/ChatGPT/jueces y controles/REVISION_ARGUMENTAL_APERTURAS_CIERRES_20260915/II/manuscript_es/sections/incidencia_icosaedrica_rev07.tex
% FORMALIZACION_NUEVA de la incidencia firmada de la carta Paley preexistente.
% La construcción comienza en exc:aw; no identifica por decreto sus banderas
% con el cociente de todas las historias TPK.
\paragraph{Construction du porteur à partir de l'incidence à six positions.}
\label{ii:ii:ico:construccion-firmada}
La matrice \(A_W\) de \eqref{ii:exc:aw}, qui intervient préalablement dans le
transport angulaire de la section~\ref{ii:sec:ii-enlace-angular-witt}, fournit
un antécédent explicite de la réalisation icosaédrique. Soit
\(C=\operatorname{bal}(A_W)\), où
\(\operatorname{bal}(0)=0\), \(\operatorname{bal}(1)=1\) et
\(\operatorname{bal}(2)=-1\). L'ordre de ses coordonnées est
\(\mathcal U=(\infty,0,1,2,3,4)\). La carte entière satisfait
\(C=C^{\mathsf T}\), \(C^2=5I_6\) et \(\operatorname{Tr}C=0\).
Nous définissons son revêtement signé par
\begin{equation}
 \mathcal D_C=\mathcal U\times\{+1,-1\},\qquad
 (u,\sigma)\sim_C(v,\tau)
 \ \Longleftrightarrow\ u\ne v\ \text{ et }\ \sigma\tau=C_{uv}.
 \label{ii:ii:ico:incidencia-firmada}
\end{equation}
L'orientation dédouble chaque coordonnée ; la matrice décide des incidences
entre les deux fibres correspondantes. L'involution de ce revêtement
est \(\jmath_C(u,\sigma)=(u,-\sigma)\).

\begin{proposition}[Réalisation isométrique de l'incidence signée]
\label{ii:ii:ico:realizacion-firmada}
L'opérateur et les vecteurs
\begin{equation}
 E_C=\frac12\left(I_6+\frac C{\sqrt5}\right),\qquad
 \mathcal H_C=\operatorname{im}E_C,\qquad
 v_{u,\sigma}=\sigma\sqrt2 E_Ce_u
 \label{ii:ii:ico:proyector-seis}
\end{equation}
réalisent \(\mathcal D_C\) comme douze vecteurs unitaires distincts dans un
espace euclidien de dimension trois. Il existe une isométrie explicite
\(T_C:\mathcal H_C\to\mathbb R^3\) pour laquelle
\(\iota_C(u,\sigma)=T_Cv_{u,\sigma}\) est une bijection sur
\(\Omega_{12}^{\rm ico}\), conserve l'incidence et satisfait
\(\iota_C\jmath_C=-\iota_C\).
\end{proposition}
\begin{proof}
De \(C^2=5I_6\), on obtient \(E_C^2=E_C=E_C^*\) ; sa trace est trois.
Les identités métriques sont
\[
 \langle v_{u,\sigma},v_{v,\tau}\rangle
 =\begin{cases}\sigma\tau,&u=v,\\
 \sigma\tau C_{uv}/\sqrt5,&u\ne v.
 \end{cases}
\]
Ainsi, les douze vecteurs sont distincts et leur incidence de produit scalaire
\(1/\sqrt5\) est exactement~\eqref{ii:ii:ico:incidencia-firmada}.
Pour expliciter l'isométrie, prenons pour colonnes de \(B\), dans l'ordre
\(\mathcal U\),
\[
 B=\frac1\nu
 \begin{pmatrix}
 0&-1&-\varphi_g&0&\varphi_g&1\\
 -1&-\varphi_g&0&1&0&-\varphi_g\\
 -\varphi_g&0&-1&-\varphi_g&-1&0
 \end{pmatrix}.
\]
La substitution de \(\varphi_g^2=\varphi_g+1\) donne
\(B^*B=I_6+C/\sqrt5=2E_C\). Ainsi
\(T_C=B/\sqrt2\big|_{\mathcal H_C}\) est une isométrie, et
\(T_Cv_{u,\sigma}=\sigma Be_u\). Les six colonnes de \(B\) et leurs
opposées sont précisément les douze sommets de
\eqref{ii:ii:pent:eq:gravity-icosahedron-vertices}. La dernière identité
démontre également la compatibilité des involutions.
\end{proof}

La réalisation dodécaphasique précédente construit les douze vecteurs
directement à partir du projecteur \(P_3^{\mathrm{ico}}\) et des positions du calendrier.
Le revêtement signé fournit une autre expression en coordonnées de ce
porteur. Pour le lecteur \(\pi_{\rm ico}\) du théorème de descente,
la relation entre les deux expressions s'écrit
\begin{equation}
 X_{\rm TPK}^{\rm exc}\xrightarrow{\ r_\pm\ }
 \mathcal D_C\xrightarrow[\simeq]{\ \iota_C\ }
 \Omega_{12}^{\rm ico},\qquad
 \pi_{\rm ico}(x)=\sigma(x)Be_{u(x)},\quad
 r_\pm(x)=(u(x),\sigma(x)).
 \label{ii:ii:ico:factorizacion-lector}
\end{equation}
Puisque \(\iota_C\) est une bijection, cette factorisation détermine
\(r_\pm=\iota_C^{-1}\pi_{\rm ico}\). Ainsi, \(r_\pm\) exprime
dans la carte signée le lecteur utilisé : son introduction
n'ajoute pas d'hypothèse indépendante pour construire le porteur
dodécaphasique. Le tableau de correspondances précédent donne cette expression
pour ses douze positions. Les propriétés (i)--(iv) conservent leur
portée propre lorsque la réalisation est identifiée au quotient
de la route exceptionnelle complète. En particulier, le transport conjoint
du repère, du projecteur et de l'incidence démontre la covariance entre cartes ;
la conservation d'une incidence fixe par une transformation active
est une identité distincte. L'orientation et la mémoire demeurent
dans l'état source pendant les deux lectures.
Une permutation signée qui transporte \(C\) transporte conjointement
\(E_C\), ses vecteurs et son incidence ; cette covariance permet de comparer
les cartes en conservant l'orientation de chaque fibre.

% END OWNER


\begin{theorem}[Descente de l’incidence au quotient icosaédrique]
\label{ii:ii:pent:thm:gravity-exceptional-descent}
Supposons que \(\pi_{\rm ico}\) satisfait (i)--(iv) et que l’action
orientée induite sur \(\Omega_{12}^{\rm ico}\) a pour groupe
\(G_{\rm ico}\simeq A_5\). Alors, pour toute classe
\(v_0\in\Omega_{12}^{\rm ico}\),
\begin{equation}
\label{ii:ii:pent:eq:gravity-icosahedral-quotient}
H_{v_0}=\operatorname{Stab}_{G_{\rm ico}}(v_0)\simeq C_5,
\qquad
\Omega_{12}^{\rm ico}\simeq G_{\rm ico}/H_{v_0}\simeq A_5/C_5.
\end{equation}
L’application
\begin{equation}
\label{ii:ii:pent:eq:gravity-coset-map}
\kappa_{v_0}:G_{\rm ico}/H_{v_0}\longrightarrow\Omega_{12}^{\rm ico},
\qquad gH_{v_0}\longmapsto g\cdot v_0,
\end{equation}
est une bijection équivariante et préserve l’incidence si l’on déclare sur le quotient
\begin{equation}
\label{ii:ii:pent:eq:gravity-coset-incidence}
gH_{v_0}\sim g'H_{v_0}
\quad\Longleftrightarrow\quad
(g\cdot v_0,g'\cdot v_0)\in\mathcal I_{12}^{\rm ico}.
\end{equation}
Par conséquent, la composée \(\kappa_{v_0}^{-1}\pi_{\rm ico}\) est l’application
typée de l’antécédent TPK vers le quotient \(A_5/C_5\).
\end{theorem}

\begin{proof}
La liste de \eqref{ii:ii:pent:eq:gravity-icosahedron-vertices} a douze éléments.
Par transitivité, \(|G_{\rm ico}\cdot v_0|=12\), et le théorème
orbite--stabilisateur donne \(|H_{v_0}|=60/12=5\). Tout groupe d’ordre
premier cinq est cyclique ; donc \(H_{v_0}\simeq C_5\). La formule de
\(\kappa_{v_0}\) ne dépend pas du représentant : si
\(gH_{v_0}=g'H_{v_0}\), alors \(g^{-1}g'\in H_{v_0}\) et
\(g'v_0=gv_0\). La surjectivité découle de la transitivité, et
l’égalité des cardinaux donne l’injectivité. L’équivariance est immédiate :
\(\kappa_{v_0}(agH_{v_0})=a\kappa_{v_0}(gH_{v_0})\). Enfin,
\eqref{ii:ii:pent:eq:gravity-coset-incidence} est précisément le transport de
structure par cette bijection ; l’hypothèse (iii) fait commuter ce
transport avec \(\pi_{\rm ico}\).
\end{proof}

Le théorème établit la descente une fois vérifiées les quatre
propriétés de \(\pi_{\rm ico}\). Son domaine est l’état enrichi
de la lecture exceptionnelle, qui conserve l’incidence transportée par
APP--TRIT--TPK. Dans cette construction, l’antécédent explicite est
l’application~\eqref{ii:ii:pent:eq:gravity-exceptional-reader-map} avec ses quatre
propriétés ; le segment allant de \(A_5/C_5\) jusqu’à la réduction transversale sans
trace est démontré ci-après avec cet antécédent et l’action induite
de \(A_5\).

La chaîne de provenance s’exprime au moyen de tous ses morphismes :
\begin{equation}
\label{ii:ii:pent:eq:gravity-exceptional-provenance}
\APP\longrightarrow\mathrm{TRIT}\longrightarrow\TPK
\longrightarrow X_{\rm TPK}^{\rm exc}
\xrightarrow{\ \pi_{\rm ico}\ }\Omega_{12}^{\rm ico}
\xleftarrow[\simeq]{\ \kappa_{v_0}\ }A_5/C_5.
\end{equation}
La composition~\eqref{ii:ii:pent:eq:gravity-exceptional-provenance} conserve
comme antécédent l’existence de \(\pi_{\rm ico}\) avec les propriétés
(i)--(iv). Sur ce domaine, la descente au quotient est le
théorème~\ref{ii:ii:pent:thm:gravity-exceptional-descent} ; on en construit
le projecteur \(P^{\mathrm{ico}}_5\), l’entrelaceur \(\Gamma_5\) et la
projection \(\Lambda(n)\). La réalisation du porteur comme champ
gravitationnel utilise, en outre, la dynamique et les contraintes physiques
spécifiées à la fin de ce développement.

Soit \(\mathcal G=A_5\), soit \(H\simeq C_5\) le stabilisateur d’un sommet
orienté de l’icosaèdre et considérons
\begin{equation}
\label{ii:ii:pent:eq:gravity-v12}
V_{12}=\R[\mathcal G/H],
\qquad |\mathcal G/H|=\frac{60}{5}=12.
\end{equation}

% BEGIN OWNER /Users/ruben/Documents/ChatGPT/jueces y controles/REVISION_ARGUMENTAL_APERTURAS_CIERRES_20260915/II/manuscript_es/sections/caracteres_pentadicos_rev07.tex
% FORMALIZACION_NUEVA: prueba explícita de la descomposición y la
% irreducibilidad ya enunciadas en el desarrollo pentacomponente de REV06.
% CERTIFICADO_NUEVO: verificar_caracteres_pentadicos_rev07.py.
% Inserción prevista tras la definición de V12, antes de su descomposición.
\paragraph{Caractères de l’incidence icosaédrique.}
Le calcul suivant s’effectue sur l’action icosaédrique obtenue dans le
théorème~\ref{ii:ii:pent:thm:gravity-exceptional-descent}, avec ses hypothèses
d’incidence et d’orientation. Il détermine la décomposition de la
représentation des sommets et l’irréductibilité de l’espace tensoriel
au moyen de leurs traces, en conservant l’antécédent
\(\pi_{\mathrm{ico}}\) déclaré dans ce théorème.

\begin{proposition}[Décomposition des sommets et caractère tensoriel]
\label{ii:prop:ii-pent-caracteres}
Soient \(\tau=(1+\sqrt5)/2\) et
\(\bar\tau=(1-\sqrt5)/2\). Dans l’ordre des classes
\(1,2A,3A,5A,5B\), les caractères de la représentation rotationnelle
\(V_3\), de sa conjuguée algébrique \(V_{3'}\), de la représentation
des douze sommets \(V_{12}\) et de
\(\operatorname{STF}_3=\operatorname{Sym}^2_0(V_3)\) sont
\begin{equation}
\label{ii:eq:ii-pent-tabla-caracteres}
\begin{array}{c|rrrrr}
\text{classe} &1&2A&3A&5A&5B\\ \hline
\text{cardinal}&1&15&20&12&12\\
\text{classe de }g^2&1&1&3A&5B&5A\\ \hline
\chi_3&3&-1&0&\tau&\bar\tau\\
\chi_{3'}&3&-1&0&\bar\tau&\tau\\
\chi_{12}&12&0&0&2&2\\
\chi_{\mathrm{STF}}&5&1&-1&0&0
\end{array}
\end{equation}
Les représentations \(V_3,V_{3'}\) et
\(\operatorname{STF}_3\) sont absolument irréductibles et distinctes.
De plus,
\begin{equation}
\label{ii:eq:ii-pent-descomposicion-demostrada}
 V_{12}\simeq\mathbf1\oplus V_3\oplus V_{3'}
                  \oplus\operatorname{STF}_3,
 \qquad
 \langle\chi_{12},\chi_{\mathrm{STF}}\rangle=1.
\end{equation}
En particulier, le facteur direct de dimension cinq apparaît une seule fois.
\end{proposition}

\begin{proof}
Dans \(A_5\), les doubles transpositions sont \(5\cdot3=15\) et
les cycles de longueur trois sont \(\binom53\cdot2=20\).
Les \(4!=24\) cycles de longueur cinq se séparent en deux classes de
douze : leur centralisateur est le sous-groupe cyclique d’ordre cinq et est
contenu dans \(A_5\). Si l’on identifie les positions d’un cycle
à \(\mathbb Z/5\mathbb Z\), la permutation qui induit
\(r\mapsto2r\) est un cycle de longueur quatre sur les éléments
non nuls, de signe négatif. Par conséquent, le carré échange les deux
classes. L’inversion \(r\mapsto-r\) est le produit de deux
transpositions et conserve chaque classe. Cela détermine les deux premières
lignes de~\eqref{ii:eq:ii-pent-tabla-caracteres}.

La représentation \(V_3\) est l’action par rotations sur les
coordonnées des sommets. Une rotation d’angle \(\theta\)
a pour trace \(1+2\cos\theta\). Les classes d’ordres deux et trois
donnent respectivement \(-1\) et \(0\). On appelle \(5A\) la
classe d’angles \(\pm2\pi/5\) et \(5B\) celle de
\(\pm4\pi/5\) ; leurs traces sont \(\tau\) et \(\bar\tau\).
Ces expressions s’obtiennent à partir de
\(u^2+u-1=0\), pour \(u=2\cos(2\pi/5)\), et de la formule
de l’angle double. L’équation de \(u\) résulte de la division de
\(1+z+z^2+z^3+z^4=0\) par \(z^2\), avec
\(z=\exp(2\pi i/5)\). On a ainsi calculé le caractère d’une
représentation déjà définie géométriquement.

Pour construire l’autre représentation, on utilise les sommets sans la
normalisation commune \(\nu^{-1}\). Leurs coordonnées appartiennent à
\(K=\mathbb Q(\sqrt5)\). Si \(B\) est la matrice de trois sommets
linéairement indépendants et que \(B_g\) contient leurs images, la
matrice rotationnelle est \(R_g=B_gB^{-1}\in M_3(K)\).
La conjugaison \(\sigma:\sqrt5\mapsto-\sqrt5\) définit alors
\(R'_g=\sigma(R_g)\), et satisfait
\[
 R'_{gh}=R'_gR'_h,\qquad
 (R'_g)^TR'_g=I,\qquad \det R'_g=1.
\]
Par conséquent, \(V_{3'}\) est une représentation réelle effectivement
construite. Son caractère est \(\sigma(\chi_3)\), qui échange
\(\tau\) et \(\bar\tau\).

Le caractère de permutation compte les sommets fixes. L’identité en fixe
douze ; les rotations non triviales d’ordre cinq fixent les deux extrémités
de leur axe. Celles d’ordre deux ou trois ne fixent aucun sommet, car le
stabilisateur d’un sommet est \(C_5\). On obtient ainsi
\(\chi_{12}=(12,0,0,2,2)\).

Pour l’espace tensoriel, si \(\lambda_1,\lambda_2,\lambda_3\)
sont les valeurs propres de \(R_g\), la somme
\(\sum_{i\leq j}\lambda_i\lambda_j\) est
\(\bigl((\sum_i\lambda_i)^2+\sum_i\lambda_i^2\bigr)/2\).
La forme métrique engendre une droite invariante dans
\(\operatorname{Sym}^2(V_3)\), dont le complément sans trace a
pour caractère
\begin{equation}
\label{ii:eq:ii-pent-caracter-stf}
 \chi_{\mathrm{STF}}(g)
 =\frac{\chi_3(g)^2+\chi_3(g^2)}2-1.
\end{equation}
En utilisant l’application de passage au carré du tableau et
\(\tau+\bar\tau=1\), \(\tau\bar\tau=-1\), on obtient
\(\chi_{\mathrm{STF}}=(5,1,-1,0,0)\).

Le produit des caractères réels est
\(\langle\chi,\psi\rangle
=\frac1{60}\sum_C|C|\chi(C)\psi(C)\).
En particulier,
\begin{align*}
 \langle\chi_3,\chi_3\rangle
 &=\langle\chi_{3'},\chi_{3'}\rangle
 =\frac{9+15+12(\tau^2+\bar\tau^2)}{60}=1,\\
 \langle\chi_3,\chi_{3'}\rangle
 &=\frac{9+15+24\tau\bar\tau}{60}=0,\\
 \langle\chi_{\mathrm{STF}},\chi_{\mathrm{STF}}\rangle
 &=\frac{25+15+20}{60}=1.
\end{align*}
La norme un prouve l’irréductibilité des complexifications ; par
conséquent, également l’irréductibilité absolue des représentations
réelles construites. Leurs dimensions et l’orthogonalité distinguent les
trois facteurs directs. Enfin, l’égalité des fonctions centrales
\[
 \chi_{12}=1+\chi_3+\chi_{3'}+\chi_{\mathrm{STF}}
\]
se vérifie dans les cinq colonnes. La semi-simplicité des
représentations d’un groupe fini en caractéristique zéro donne la
décomposition~\eqref{ii:eq:ii-pent-descomposicion-demostrada}. La
multiplicité du facteur direct tensoriel peut se vérifier directement :
\[
 \langle\chi_{12},\chi_{\mathrm{STF}}\rangle
 =\frac{12\cdot5}{60}=1.
\]
\end{proof}

Ainsi, l’idempotent central associé à
\(\chi_{\mathrm{STF}}=\chi_5\) a pour rang
\(5\cdot1=5\) dans \(V_{12}\). L’irréductibilité qui intervient
dans la normalisation de l’entrelaceur tensoriel et la multiplicité qui
intervient dans le couplage gravitationnel sont attestées par
ce même calcul.

% END OWNER


La représentation de permutation \(\rho_{12}\) se décompose comme
\begin{equation}
\label{ii:ii:pent:eq:gravity-v12-decomposition}
V_{12}\simeq\mathbf1\oplus V_3\oplus V_{3'}\oplus V_5.
\end{equation}
L’idempotent central du facteur de caractère
\(\chi_5=(5,1,-1,0,0)\) est
\begin{equation}
\label{ii:ii:pent:eq:gravity-p5}
\boxed{
P^{\mathrm{ico}}_5=\frac5{60}\sum_{g\in A_5}\chi_5(g^{-1})\rho_{12}(g)
=\frac1{12}\left(
5I+\sum_{g\in2A}\rho_{12}(g)
-\sum_{g\in3A}\rho_{12}(g)
\right).}
\end{equation}

\begin{theorem}[Projecteur et couplage gravitationnel de rang cinq]
\label{ii:ii:pent:thm:gravity-p5-coupling}
L'opérateur \(P^{\mathrm{ico}}_5\) satisfait
\begin{equation}
\label{ii:ii:pent:eq:gravity-p5-properties}
(P^{\mathrm{ico}}_5)^2=P^{\mathrm{ico}}_5=(P^{\mathrm{ico}}_5)^*,
\qquad \Tr P^{\mathrm{ico}}_5=\operatorname{rank} P^{\mathrm{ico}}_5=5.
\end{equation}
Pour
\begin{equation}
\label{ii:ii:pent:eq:gravity-G5}
\mathbb G_{5,a}:=G_a(\vartheta^\circ_{108})P^{\mathrm{ico}}_5
\end{equation}
on a
\begin{equation}
\label{ii:ii:pent:eq:gravity-G5-spectrum}
\Spec(\mathbb G_{5,a})
=\{G_a(\vartheta^\circ_{108})\text{ avec multiplicité }5,
\;0\text{ avec multiplicité }7\}.
\end{equation}
\end{theorem}

\begin{proof}
L’orthogonalité des caractères fait de \eqref{ii:ii:pent:eq:gravity-p5}
l’idempotent central qui agit comme l’identité sur l’unique copie de
\(V_5\) et comme zéro sur les trois autres facteurs de
\eqref{ii:ii:pent:eq:gravity-v12-decomposition}. La réalité du caractère et
l’orthogonalité de \(\rho_{12}\) donnent l’autoadjonction. Sa trace et son rang
sont cinq. Multiplier un projecteur de rang cinq par le scalaire positif
\(G_a\) produit immédiatement le spectre indiqué.
\end{proof}

La section gravitationnelle \(G_a\) est commune aux cinq composantes du
facteur direct tensoriel sélectionné. Sa multiplicité spectrale exprime le
rang du porteur ; les sept autres modes du module de permutation
appartiennent au noyau de \(\mathbb G_{5,a}\). Le symbole
\(P^{\mathrm{ico}}_5\) identifie ici ce projecteur central de rang
cinq sur \(V_{12}\) ; la moyenne des cinq positions d’une
pentafibre arithmétique a un autre domaine et est de rang un.

\subsection{Entrelaceur avec l’espace des tenseurs symétriques sans trace}

Soit \(\mathcal V=\Omega_{12}^{\rm ico}\subset S^2\) l’ensemble des
douze sommets unitaires déjà construit. Pour chaque \(v\in\mathcal V\), définissons
\begin{equation}
\label{ii:ii:pent:eq:gravity-Qv}
Q_v=vv^T-\frac13I_3\in
\operatorname{STF}_3:=\operatorname{Sym}^2_0(\R^3).
\end{equation}
L’application de repère
\begin{equation}
\label{ii:ii:pent:eq:gravity-Gamma0}
\Gamma_0\!\left(\sum_{v\in\mathcal V}x_ve_v\right)
=\sum_{v\in\mathcal V}x_vQ_v
\end{equation}
est \(A_5\)-équivariante et satisfait
\begin{equation}
\label{ii:ii:pent:eq:gravity-tight-frame}
\Gamma_0\Gamma_0^*=\frac85I_{\operatorname{STF}_3},
\qquad
\Gamma_0=\Gamma_0P^{\mathrm{ico}}_5.
\end{equation}

\begin{theorem}[Entrelaceur tensoriel normalisé]
\label{ii:ii:pent:thm:gravity-gamma5}
L’application
\begin{equation}
\label{ii:ii:pent:eq:gravity-gamma5}
\boxed{
\Gamma_5:=\sqrt{\frac58}\,\Gamma_0|_{V_5}:
V_5\xrightarrow{\ \simeq\ }\operatorname{STF}_3}
\end{equation}
est une isométrie \(A_5\)-équivariante. Son adjoint est
\begin{equation}
\label{ii:ii:pent:eq:gravity-gamma5-adjoint}
\Gamma_5^*Q=\sqrt{\frac58}
\sum_{v\in\mathcal V}(v^TQv)e_v.
\end{equation}
\end{theorem}

\begin{proof}
L’équivariance découle de \(Q_{gv}=gQ_vg^T\). L’opérateur
\(\Gamma_0\Gamma_0^*\) commute avec la représentation irréductible de
\(A_5\) sur \(\operatorname{STF}_3\) ; il est donc scalaire. Sa trace
est
\[
\sum_{v\in\mathcal V}\|Q_v\|_F^2
=12\cdot\frac23=8;
\]
en divisant par la dimension cinq, on obtient \(8/5\). D’où
\(\Gamma_5\Gamma_5^*=I\). Comme le domaine et le codomaine sont de dimension
cinq, on a aussi \(\Gamma_5^*\Gamma_5=I\). La formule de l’adjoint
découle de \(\langle Q,Q_v\rangle_F=v^TQv\).
\end{proof}

La normalisation par le repère icosaédrique construit l’isométrie qui
transporte le facteur fini vers l’espace tensoriel réel de spin deux.

\subsection{Projection transversale sans trace}

Pour \(n\in S^2\), soit \(\Pi(n)=I-nn^T\). Nous définissons
\begin{equation}
\label{ii:ii:pent:eq:gravity-lambda-tt}
\boxed{
\Lambda(n)Q
=\Pi Q\Pi-\frac12\Tr(\Pi Q\Pi)\Pi.}
\end{equation}

\begin{theorem}[Réduction de cinq composantes à deux]
\label{ii:ii:pent:thm:gravity-tt}
La restriction de \(\Lambda(n)\) à \(\operatorname{STF}_3\) est le
projecteur orthogonal sur
\begin{equation}
\label{ii:ii:pent:eq:gravity-htt}
\mathcal H_{\rm TT}(n)
=\{Q\in\operatorname{STF}_3:Qn=0\}.
\end{equation}
En particulier,
\begin{equation}
\label{ii:ii:pent:eq:gravity-lambda-properties}
\Lambda(n)^2=\Lambda(n)=\Lambda(n)^*,
\qquad \operatorname{rank}\Lambda(n)=2.
\end{equation}
Pour \(n=e_3\), l’image est engendrée par
\begin{equation}
\label{ii:ii:pent:eq:gravity-plus-cross}
E_+=\frac1{\sqrt2}\operatorname{diag}(1,-1,0),
\qquad
E_\times=\frac1{\sqrt2}
\begin{pmatrix}0&1&0\\1&0&0\\0&0&0\end{pmatrix}.
\end{equation}
\end{theorem}

\begin{proof}
Comme \(\Pi n=0\), la sortie est transversale. Sa trace vaut
\[
\Tr(\Pi Q\Pi)-\frac12\Tr(\Pi)\Tr(\Pi Q\Pi)=0,
\]
car \(\Tr\Pi=2\). Si \(Qn=0\) et \(\Tr Q=0\), alors
\(\Pi Q\Pi=Q\) et \(\Tr(\Pi Q\Pi)=0\) ; donc
\(\Lambda(n)Q=Q\). La formule est autoadjointe pour le produit de
Frobenius. Dans le repère \(n=e_3\), sa matrice dans une base STF adaptée est
\(\operatorname{diag}(1,1,0,0,0)\), ce qui prouve le rang.
\end{proof}

La chaîne complète devient donc
\begin{equation}
\label{ii:ii:pent:eq:gravity-12-5-5-2}
\boxed{
V_{12}\xrightarrow{P^{\mathrm{ico}}_5}V_5
\xrightarrow{\Gamma_5}\operatorname{STF}_3
\xrightarrow{\Lambda(n)}\mathcal H_{\rm TT}(n),
\qquad 12\longrightarrow5\longrightarrow5\longrightarrow2.}
\end{equation}
Le couplage réduit
\begin{equation}
\label{ii:ii:pent:eq:gravity-effective-map}
\mathbb G_{{\rm TT},a}(n)
=G_a(\vartheta^\circ_{108})\Lambda(n)\Gamma_5P^{\mathrm{ico}}_5
\end{equation}
est de rang deux.

\subsection{Décomposition en \(5\) modes de polarisation}

La dimension cinq acquiert un contenu tensoriel lorsque l’on choisit une direction
unitaire de propagation \(n\) et un repère orthonormé orienté
\((e_1,e_2,n)\). Nous définissons les cinq tenseurs réels
\begin{align}
E_{+}&=\frac1{\sqrt2}(e_1e_1^T-e_2e_2^T),
&E_{\times}&=\frac1{\sqrt2}(e_1e_2^T+e_2e_1^T),
\label{ii:ii:pent:eq:gravity-helicity2-real}\\
E_{1}&=\frac1{\sqrt2}(e_1n^T+ne_1^T),
&E_{2}&=\frac1{\sqrt2}(e_2n^T+ne_2^T),
\label{ii:ii:pent:eq:gravity-helicity1-real}\\
E_{0}&=\frac1{\sqrt6}(2nn^T-e_1e_1^T-e_2e_2^T).
\label{ii:ii:pent:eq:gravity-helicity0-real}
\end{align}
Ils sont tous symétriques et sans trace. Les deux premiers sont transversaux ; les
deux suivants mêlent le plan transversal à la direction de
propagation, et le dernier est le mode scalaire longitudinal sans trace.

\begin{theorem}[Décomposition hélicoïdale complète du porteur]
\label{ii:ii:pent:thm:gravity-five-helicities}
Les tenseurs de
\eqref{ii:ii:pent:eq:gravity-helicity2-real}, \eqref{ii:ii:pent:eq:gravity-helicity1-real}, \eqref{ii:ii:pent:eq:gravity-helicity0-real}
forment une base orthonormée de \(\operatorname{STF}_3\). Dans la
complexification, les combinaisons
\begin{equation}
\label{ii:ii:pent:eq:gravity-complex-helicities}
E_{\pm2}=\frac1{\sqrt2}(E_+\mp i E_\times),
\qquad
E_{\pm1}=\frac1{\sqrt2}(E_1\mp i E_2),
\qquad E_0=E_0
\end{equation}
sont des vecteurs de poids \(\pm2,\pm1,0\), respectivement, sous les rotations
autour de \(n\). De plus,
\begin{equation}
\label{ii:ii:pent:eq:gravity-tt-on-five}
\Lambda(n)E_{\pm2}=E_{\pm2},
\qquad
\Lambda(n)E_{\pm1}=0,
\qquad
\Lambda(n)E_0=0.
\end{equation}
\end{theorem}

\begin{proof}
Dans le repère \((e_1,e_2,n)\), chaque tenseur a une norme de Frobenius égale à un.
Leurs motifs d’entrées non diagonales sont disjoints, et les trois diagonales
de \(E_+,E_0\) ont un produit intérieur nul ; les cinq tenseurs
sont donc orthonormés. Comme \(\dim\operatorname{STF}_3=6-1=5\), ils forment une base.

Une rotation d’angle \(\psi\) autour de \(n\) envoie
\((e_1,e_2)\) sur
\((\cos\psi\,e_1+\sin\psi\,e_2,
-\sin\psi\,e_1+\cos\psi\,e_2)\). En substituant dans les formules, on obtient
une rotation d’angle \(2\psi\) sur
\(\operatorname{span}\{E_+,E_\times\}\), une d’angle \(\psi\) sur
\(\operatorname{span}\{E_1,E_2\}\), et \(E_0\) fixe. D’où les poids de
\eqref{ii:ii:pent:eq:gravity-complex-helicities}. Enfin,
\(\Pi E_{+}\Pi=E_+\) et \(\Pi E_\times\Pi=E_\times\), tandis que
\(\Pi E_1\Pi=\Pi E_2\Pi=0\). Pour \(E_0\),
\(\Pi E_0\Pi=-(e_1e_1^T+e_2e_2^T)/\sqrt6\), qui est purement de trace dans
le plan, et la soustraction de \eqref{ii:ii:pent:eq:gravity-lambda-tt} l’annule.
\end{proof}

Le théorème précise le sens de ``cinq polarisations'' sans leur attribuer
à toutes le même statut dynamique. Avant d’imposer les équations de champ,
le porteur de spin deux a cinq coordonnées irréductibles. Dans une
réalisation massive de Fierz--Pauli, les contraintes laissent précisément
les cinq poids de \eqref{ii:ii:pent:eq:gravity-complex-helicities} ; dans une réalisation
sans masse avec invariance de jauge, les secteurs \(\pm1\) et \(0\) sont éliminés et
il reste \(\pm2\). HMT fournit ici le porteur, la base et le projecteur ;
le choix de la dynamique massive ou sans masse relève de l’interface
physique. Cinq ne signifie donc pas ``cinq ondes observées'', mais cinq
degrés tensoriels disponibles avant la réduction.

L’orientation reste également visible : remplacer
\((e_1,e_2,n)\) par \((e_1,-e_2,-n)\) conserve \(E_+,E_0\) et échange
les phases des paires hélicoïdales. Cette information ne se trouve pas dans le
scalaire \(G_a\) ; elle réside dans le porteur sur lequel agit
\(\mathbb G_{5,a}\).

\subsection{Interprétation physique de la chaîne représentationnelle}

La dimension cinq apparaît dans des objets distincts, et seul
\eqref{ii:ii:pent:eq:gravity-12-5-5-2} autorise le transport entre eux :

\begin{enumerate}
\item \(V_5\) est le facteur irréductible de la représentation sur les
douze sommets. C’est un résultat fini de théorie des représentations.
\item \(\operatorname{STF}_3\) est le porteur réel de dimension cinq des
poids \(m=-2,-1,0,1,2\) de spin deux.
\item Une théorie massive de Fierz--Pauli peut réaliser ces cinq poids
comme cinq états physiques, après fixation de l’action et des contraintes.
\item La relativité générale sans masse conserve deux hélicités radiatives,
\(\pm2\), représentées par \(E_+\) et \(E_\times\).
\item Les classes \(E(2)\) d’ondes métriques, comme \(III_5\), classent
les amplitudes possibles sous d’autres hypothèses ; elles ne sont pas automatiquement
\(V_5\), Fierz--Pauli ou la relativité générale.
\end{enumerate}

Le contenu structurel est ainsi déterminé : sous les conditions de la
descente d’incidence, HMT construit un porteur tensoriel à cinq
composantes avant les contraintes. La réalisation sans masse conserve
les deux hélicités radiatives au moyen de la projection TT de rang deux.
L’information d’orientation se conserve dans les tenseurs et dans leur
transport ; la valeur scalaire \(G_a\) détermine le couplage commun.
